\chapter{问题定义}\label{ch:problem-definition}

\section{任务边界}

本系统支持的核心任务是：给定医学研究对象和数据需求，发现并采集一个或多个科学数据来源，将其转换为统一的单表或多表结构，保留来源和处理记录，并输出便于后续分析的CSV 数据产品。

相应地，本系统不把下列事项作为正式交付承诺：自动提出并验证科学假设，替代领域专家判断数据的科学含义，对研究结果给出医学结论，保证互联网范围内绝对查全，或将模型抽取值自动升级为高可信事实。

Core 按执行规格检查候选结果。必填字段、单位、外键或必需表不满足规则时，系统停止发布并记录原因。允许缺失的可选字段可以保留为空，但必须说明缺口。发布成功因此只表示本次文件通过了配置的检查，不表示原始科研需求全部完成。第\ref{ch:experiments}章样例10需要四张关联表，其中必需数据未能取得，系统没有用占位值补齐，而是以“未发布阻塞”结束。

\section{设计目标}\label{sec:design-goals}

\begin{enumerate}
  \item 可用：输出具有明确字段和表关系的 CSV，供 R、Python 或其他统计软件读取。
  \item 可追溯：登记来源身份与文件摘要，并尽可能记录页码、表格行列、数据库记录号等定位信息。
  \item 可检查：由注册代码执行解析、规范化、整合和发布；缺失、重复、冲突及单位问题进入检查记录。
  \item 可修正：无法自动判断时请求结构化反馈，保存原值、决定与修改结果。
  \item 可恢复：保存运行事件与下载进度，为重试、断点续传和审批后恢复提供依据；不能据此承诺任意步骤均无需重做。
  \item 如实交付：明确区分工作区文件、候选结果与正式发布，同时说明原始需求中尚未完成的部分。
  \item 可扩展：高频场景使用预先注册的数据族；未注册的多表任务可以提交动态规格。数据族是表结构、解析方法与检查规则的注册单位。
\end{enumerate}

系统重点防范以下失败：来源URL 或版本漂移，下载中断或内容类型错误，字段映射模糊，单位、尺度或行粒度不兼容，重复与冲突被静默覆盖，模型输出缺少证据，任务重启后重复或改变处理，过期执行晚到发布，以及正式文件被修改后仍被展示。这些失败模式对应的机制与实测拦截记录一览见\ref{sec:reliability-base} 节。

系统自动处理规则能够明确判断的问题；其余问题应停止处理、保留说明或请求人工判断，而不是猜测后继续发布。

